\documentclass[10pt]{article}
\usepackage[utf8]{inputenc}
\usepackage[T1]{fontenc}
\usepackage[spanish,es-nodecimaldot]{babel}
\usepackage[margin=1.4cm]{geometry}
\usepackage{microtype}
\setlength{\parindent}{0pt}
\setlength{\parskip}{2pt}
\pagestyle{empty}

\begin{document}
\begin{center}
{\large\bfseries Torneo de Komikan concurrente}\\[-1pt]
{\small Tarea 2 --- Sistemas Operativos}
\end{center}
\vspace{-5pt}

\section*{1. Diseño del Esquema de Sincronización}
Cada mesa ejecuta trece hilos de fichas (Puma y doce Perros) y un hilo Juez.
El estado se divide por responsabilidad: \texttt{tablero.c} implementa el
monitor \texttt{Casillas}, mientras \texttt{mesa.c} aplica las reglas y
coordina el estado lógico de la partida.

El tablero tiene 28 nodos. Cada nodo es un recurso independiente con
\texttt{owner[c]}, protegido al consultarlo o modificarlo por su propio
\texttt{mtx[c]}; \texttt{cond[c]} permite esperar hasta que quede libre.
Importante: ese mutex se mantiene solo durante la operación breve sobre
\texttt{owner}; la exclusión del tablero es lógica y persiste mientras la ficha
sea su dueña, no consiste en retener un mutex POSIX durante el movimiento.
Así, fichas con regiones disjuntas pueden operar simultáneamente. El Perro
reserva su origen y vecinos (hasta 5 nodos); el Puma agrega los destinos de
salto (hasta 9).

Antes de leer o cambiar \texttt{tablero[]}, cada ficha construye y ordena
ascendentemente su vecindario, y reserva todos sus nodos. La validación exige
que origen, destino y, para una captura, casilla saltada estén reservados.
Entonces \texttt{juez\_validar} consulta el estado autoritativo y
\texttt{juez\_aplicar} actualiza el tablero; la ficha conserva las reservas
hasta registrar el movimiento y liberar la región en orden inverso.
\texttt{vista[]} es una copia informativa: \texttt{mtx\_mesa} protege la
instantánea, estados, contadores y registro, pero no sustituye los monitores
que protegen \texttt{tablero[]}. La jerarquía es casillas (por índice
ascendente), \texttt{mtx\_mesa} y, al emitir mensajes, consola. El Juez no toma
casillas. Los datos se inicializan antes de crear los hilos y se consultan al
final después de unirlos; \texttt{terminar} y los latidos son atómicos.

\section*{2. Análisis de Riesgos de Concurrencia}
\textbf{Deadlock.} La condición necesaria de espera circular se elimina
imponiendo un orden total: todos reservan nodos en orden ascendente y liberan
en orden inverso. Si una ficha espera un nodo, ya posee solo nodos menores;
por tanto, las aristas de cualquier cadena de espera avanzan estrictamente
por índice y no pueden cerrar un ciclo. Además, \texttt{tomar} libera el
mutex POSIX antes de esperar en la condición, y si la partida termina durante
la reserva, \texttt{casillas\_bloquear} deshace las reservas parciales.
El Juez nunca espera por una casilla mientras mantiene \texttt{mtx\_mesa}.
La ruta de fallo simulado también se recupera: el Juez detecta la falta de
latido, une los hilos y libera las reservas huérfanas.

\textbf{Livelock.} La contención no se resuelve con reintentos activos:
quien encuentra una casilla ocupada espera en \texttt{cond[c]} por tramos de
50\,ms, actualiza su latido y vuelve a comprobar el predicado; al liberar una
casilla se señaliza a un esperador. Cada ficha prueba un conjunto finito de
candidatos y luego libera su región, y sus ciclos incluyen un retardo
configurable. Esto evita que el protocolo de reserva consuma CPU girando.
No obstante, no existe garantía de progreso de cada movimiento: estrategias
sin mejora, conflictos o falta de movimientos pueden impedirlo. El Juez
detecta al Puma acorralado y termina por condición de victoria, abandono o
duración máxima configurada (empate), evitando una partida indefinida.

\textbf{Starvation.} La señalización y el orden de adquisición no son una
política FIFO: POSIX no garantiza qué esperador despierta primero y tampoco
hay turnos ni cuota por ficha. Por ello, no se puede afirmar que el código
impida formalmente la inanición de una ficha bajo cualquier planificador.
La espera por condición, los ámbitos locales de reserva y los latidos
actualizados mientras se espera reducen espera activa y falsas detecciones de
fallo, pero no constituyen una garantía de equidad. El límite temporal de la
partida acota cuánto puede prolongarse el torneo; no garantiza que todos se
muevan antes del empate.

\section*{3. Evaluación}
La estrategia es sencilla de auditar: monitores por nodo evitan un mutex
global del tablero, y la validación bajo la región reservada mantiene
consistentes las reglas y las capturas. Con tablero fijo de 28 nodos, formar y
ordenar el vecindario cuesta $O(V^2)$ por inserción, con $V\leq 9$; la
deduplicación y la validación de hasta ocho propuestas también tienen cotas
constantes pequeñas. El coste real del movimiento incluye la contención y la
salida a log/consola, que puede prolongar la sección de \texttt{mtx\_mesa}.

La latencia no es de tiempo real: depende del retardo configurado, la
planificación y la ocupación de los nodos. Una liberación señaliza al
esperador; el tramo de 50\,ms sirve para reexaminar la condición, no es una
demora obligatoria para todo movimiento. El Juez reevalúa cada 100\,ms o al
recibir eventos como captura o Puma acorralado. En una mesa pueden estar
activas hasta 13 fichas, con movimientos concurrentes cuando sus vecindarios
no se solapan; con $M$ mesas, el límite es $13M$ hilos de fichas, además de
Jueces y servidor. Se gana concurrencia frente al bloqueo global, a cambio
de más primitivas y de no ofrecer equidad estricta; el mutex de mesa y la
impresión detallada son los principales puntos de serialización.

\end{document}
